\section{数据定价：数据不是按 GB 卖}

数据被称为新时代的石油，但这个比喻容易误导。石油按物理量卖，数据的价值却取决于真实性、稀缺性、可验证性、可复用性和边际收益。同样 1GB 数据，可能是重复网页，也可能是关键长尾失败案例；价值完全不同。

\begin{figure}[H]
\centering
\includegraphics[width=0.92\textwidth]{data-pricing.png}
\caption{数据定价的五个维度：真实性、稀缺性、可验证性、可复用性、边际收益。}
\end{figure}

\begin{knowledgebox}{读图：数据价格来自边际贡献}
图中五个因素共同决定数据价值。越真实、越稀缺、越可验证、越能迁移，数据越值钱；但同类数据越多，边际收益越低。数据产业的核心不是堆量，而是找到高边际收益数据。
\end{knowledgebox}

\subsection{失败数据为什么贵}

失败数据贵，因为它告诉模型边界在哪里。成功轨迹可能很多都相似，失败轨迹往往暴露系统缺陷、长尾场景和恢复策略。对自动驾驶、机器人和 Agent 来说，失败-修复对比单纯成功演示更有训练价值。

\begin{importantbox}{最有效的数据：失败再成功}
先失败再成功的数据包含错误、诊断、修正和结果。它比单纯成功数据更接近学习过程，也更适合训练恢复能力。
\end{importantbox}

\subsection{本章小结}

数据定价应按能力提升的边际贡献，而不是按体积。高价值数据通常真实、稀缺、可验证，并能暴露模型边界。

